% GENERATED FILE. Edit canonical lessons, then run npm run generate:books.
% canonical-source-hash: fnv1a64:337794bac8606b77
% canonical-lessons: TE-C300-butlu, TE-C300-golusu, TE-C300-cevipogulu, TE-C300-gundi, TE-C300-cetirumalu

\chapter{Shoes, Chain, Earrings, Button, Handkerchief}
\label{ch:te-shoes-chain-earrings-button-handkerchief}
\hlchaptermodality{300}

\begin{chapteropening}
\textbf{By the end of this chapter:} \emph{I can say shoes, a chain, earrings, a button and a handkerchief.}

Complete the last lesson of chapter 300: I can say shoes, a chain, earrings, a button and a handkerchief.
\end{chapteropening}

\section[\texorpdfstring{\te{బూట్లు} (būṭlu)}{బూట్లు (būṭlu)}]{\te{బూట్లు} (būṭlu) — shoes}

\label{lesson:TE-C300-butlu}



\begin{quote}
Before the new one: say the Telugu for curry leaves, then the Telugu for buttermilk.
\end{quote}

\subsection*{You'll want to know: \te{బూట్లు}}
\textbf{\te{బూట్లు}} — \emph{būṭlu} — \textquotedblleft{}shoes\textquotedblright{}.

\begin{grammarlens}[title={What you've built}]
One more word for this chapter.
\end{grammarlens}

\subsection*{Your turn}
\begin{itemize}
\raggedright
  \item \emph{Say it:} \emph{būṭlu}
  \item \emph{Say it:} \emph{būṭlu}, once more
  \item \emph{Say it:} say \emph{būṭlu}
  \item \emph{Recall:} say the Telugu for curry leaves, then the Telugu for buttermilk, then say \emph{būṭlu} again
\end{itemize}

\begin{tcolorbox}[breakable,colback=teal!4,colframe=teal!35!black,title={Before you move on}]
What does \te{బూట్లు} mean? (\textquotedblleft{}Shoes\textquotedblright{}.) Say it once more.
\end{tcolorbox}
\section[\texorpdfstring{\te{గొలుసు} (golusu)}{గొలుసు (golusu)}]{\te{గొలుసు} (golusu) — a chain, a necklace}

\label{lesson:TE-C300-golusu}



\begin{quote}
Before the new one: say the Telugu for buttermilk, then the Telugu for shoes.
\end{quote}

\subsection*{You'll want to know: \te{గొలుసు}}
\textbf{\te{గొలుసు}} — \emph{golusu} — \textquotedblleft{}a chain, a necklace\textquotedblright{}.

\begin{grammarlens}[title={What you've built}]
One more word for this chapter.
\end{grammarlens}

\subsection*{Your turn}
\begin{itemize}
\raggedright
  \item \emph{Say it:} \emph{golusu}
  \item \emph{Say it:} \emph{golusu}, once more
  \item \emph{Say it:} say \emph{golusu}
  \item \emph{Recall:} say the Telugu for buttermilk, then the Telugu for shoes, then say \emph{golusu} again
\end{itemize}

\begin{tcolorbox}[breakable,colback=teal!4,colframe=teal!35!black,title={Before you move on}]
What does \te{గొలుసు} mean? (\textquotedblleft{}A chain, a necklace\textquotedblright{}.) Say it once more.
\end{tcolorbox}
\section[\texorpdfstring{\te{చెవిపోగులు} (cevipōgulu)}{చెవిపోగులు (cevipōgulu)}]{\te{చెవిపోగులు} (cevipōgulu) — earrings}

\label{lesson:TE-C300-cevipogulu}



\begin{quote}
Before the new one: say the Telugu for shoes, then the Telugu for a chain.
\end{quote}

\subsection*{You'll want to know: \te{చెవిపోగులు}}
\textbf{\te{చెవిపోగులు}} — \emph{cevipōgulu} — \textquotedblleft{}earrings\textquotedblright{}.

\begin{grammarlens}[title={What you've built}]
One more word for this chapter.
\end{grammarlens}

\subsection*{Your turn}
\begin{itemize}
\raggedright
  \item \emph{Say it:} \emph{cevipōgulu}
  \item \emph{Say it:} \emph{cevipōgulu}, once more
  \item \emph{Say it:} say \emph{cevipōgulu}
  \item \emph{Recall:} say the Telugu for shoes, then the Telugu for a chain, then say \emph{cevipōgulu} again
\end{itemize}

\begin{tcolorbox}[breakable,colback=teal!4,colframe=teal!35!black,title={Before you move on}]
What does \te{చెవిపోగులు} mean? (\textquotedblleft{}Earrings\textquotedblright{}.) Say it once more.
\end{tcolorbox}
\section[\texorpdfstring{\te{గుండీ} (guṇḍī)}{గుండీ (guṇḍī)}]{\te{గుండీ} (guṇḍī) — a button}

\label{lesson:TE-C300-gundi}



\begin{quote}
Before the new one: say the Telugu for a chain, then the Telugu for earrings.
\end{quote}

\subsection*{You'll want to know: \te{గుండీ}}
\textbf{\te{గుండీ}} — \emph{guṇḍī} — \textquotedblleft{}a button\textquotedblright{}.

\begin{grammarlens}[title={What you've built}]
One more word for this chapter.
\end{grammarlens}

\subsection*{Your turn}
\begin{itemize}
\raggedright
  \item \emph{Say it:} \emph{guṇḍī}
  \item \emph{Say it:} \emph{guṇḍī}, once more
  \item \emph{Say it:} say \emph{guṇḍī}
  \item \emph{Recall:} say the Telugu for a chain, then the Telugu for earrings, then say \emph{guṇḍī} again
\end{itemize}

\begin{tcolorbox}[breakable,colback=teal!4,colframe=teal!35!black,title={Before you move on}]
What does \te{గుండీ} mean? (\textquotedblleft{}A button\textquotedblright{}.) Say it once more.
\end{tcolorbox}
\section[\texorpdfstring{\te{చేతిరుమాలు} (cētirumālu)}{చేతిరుమాలు (cētirumālu)}]{\te{చేతిరుమాలు} (cētirumālu) — a handkerchief}

\label{lesson:TE-C300-cetirumalu}



\begin{quote}
Before the new one: say the Telugu for earrings, then the Telugu for a button.
\end{quote}

\subsection*{You'll want to know: \te{చేతిరుమాలు}}
\textbf{\te{చేతిరుమాలు}} — \emph{cētirumālu} — \textquotedblleft{}a handkerchief\textquotedblright{}.

\begin{grammarlens}[title={What you've built}]
One more word for this chapter.
\end{grammarlens}

\subsection*{Your turn}
\begin{itemize}
\raggedright
  \item \emph{Say it:} \emph{cētirumālu}
  \item \emph{Say it:} \emph{cētirumālu}, once more
  \item \emph{Say it:} say \emph{cētirumālu}
  \item \emph{Recall:} say the Telugu for earrings, then the Telugu for a button, then say \emph{cētirumālu} again
\end{itemize}

\begin{tcolorbox}[breakable,colback=teal!4,colframe=teal!35!black,title={Before you move on}]
What does \te{చేతిరుమాలు} mean? (\textquotedblleft{}A handkerchief\textquotedblright{}.) Say it once more.
\end{tcolorbox}
